% =========================================================
% IEEE Conference Research Proposal
% A Machine Learning Framework for Pitting Corrosion Detection
% in Offshore Infrastructure Using Fused RGB, Thermal, and LiDAR Imagery
%
% Author: Victor — MSc by Research, Universiti Teknologi PETRONAS
%
% To compile: pdflatex proposal.tex  (run twice for cross-references)
%
% References: replace \bibitem placeholders with your actual sources.
% To switch to BibTeX, replace \begin{thebibliography}...\end{thebibliography}
% with \bibliographystyle{IEEEtran} and \bibliography{references}
% =========================================================

\documentclass[conference]{IEEEtran} 
\IEEEoverridecommandlockouts

\usepackage{cite}
\usepackage{amsmath,amssymb,amsfonts}
\usepackage{graphicx}
\usepackage{textcomp}
\usepackage{xcolor}
\usepackage{booktabs}

\begin{document}

\title{A Machine Learning Framework for Pitting Corrosion Detection in Offshore
Infrastructure Using Fused RGB, Thermal, and LiDAR Imagery}

\author{
    \IEEEauthorblockN{Victor}
    \IEEEauthorblockA{
        \textit{Department of Electrical and Electronic Engineering}\\
        \textit{Universiti Teknologi PETRONAS (UTP)}\\
        Seri Iskandar, Perak, Malaysia
    }
}

\maketitle

% ---------------------------------------------------------
\begin{abstract}
Pitting corrosion is one of the most hazardous and economically
costly forms of material degradation affecting offshore oil and gas
infrastructure, including fixed platforms, subsea pipelines, and
floating production systems. Conventional automated inspection
approaches rely predominantly on single-modality sensing—visible-spectrum
(RGB) imagery, thermal infrared, or structured-light scanning in
isolation—which constrains detection reliability under the variable
illumination, surface fouling, and harsh environmental conditions
characteristic of offshore deployments. This paper proposes a
multimodal machine learning (ML) framework that fuses RGB, thermal
infrared, and LiDAR point-cloud data for the automated detection
and localisation of pitting corrosion on offshore structures. By
exploiting the complementary strengths of each modality—colour and
texture discrimination from RGB, subsurface thermal anomaly signatures
from infrared, and high-fidelity geometric surface deviation from
LiDAR—the proposed hybrid feature-level fusion architecture aims
to achieve detection accuracy and environmental robustness that no
single-modality approach can deliver alone. The framework is to be
evaluated against single- and dual-modality baselines using
standardised metrics including intersection over union (IoU),
precision, recall, and F1-score. This research addresses a
documented gap in the literature: no existing study has unified
all three sensor modalities specifically for pitting corrosion
detection in offshore environments. Expected contributions include
a novel trimodal fusion framework, a benchmark multimodal corrosion
dataset, and quantified accuracy improvements, with direct implications
for structural safety and inspection efficiency in the offshore
energy sector.
\end{abstract}

\begin{IEEEkeywords}
pitting corrosion, multimodal fusion, LiDAR, thermal infrared imaging,
convolutional neural networks, offshore structural health monitoring
\end{IEEEkeywords}

% ---------------------------------------------------------
\section{Introduction}

Offshore oil and gas infrastructure operates in one of the most
corrosive environments on Earth. Sustained exposure to marine aerosols,
saltwater immersion, cyclic thermal loading, and biological fouling
accelerates corrosion processes that, if undetected, can precipitate
catastrophic structural failure, environmental contamination, and
substantial economic loss \cite{ref1}. Among the various corrosion
morphologies, pitting corrosion is particularly insidious: its
localised, subsurface-nucleating character means that significant
material loss can develop with minimal macroscopic surface signature
until the corrosion penetrates through the wall of a pipe or plate
\cite{ref2}.

Current offshore inspection regimes combine periodic manual visual
inspection with single-sensor non-destructive testing (NDT) techniques.
Whilst RGB camera systems have gained traction for automated surface
inspection due to their low deployment cost, they degrade markedly
under poor illumination, surface contamination, and adverse weather
\cite{ref3}. Thermal infrared (IR) imaging can reveal subsurface
thermal anomalies indicative of material thinning but is limited in
spatial resolution and prone to false positives from environmentally
induced thermal gradients \cite{ref4}. LiDAR sensors provide dense
three-dimensional point clouds encoding surface topology with
sub-millimetre precision, enabling pit detection as geometric
depressions, yet are susceptible to noise from water films, maritime
aerosols, and specular reflectance on metallic substrates \cite{ref5}.

Multimodal sensor fusion—exploiting complementary sensing physics to
improve system-level performance—is well established in autonomous
driving, aerial remote sensing, and medical imaging. However, its
systematic application to pitting corrosion detection on offshore
structures remains largely unexplored. This gap represents both a
scientific opportunity and a practical imperative given the industry's
need for more reliable, autonomous inspection solutions.

This research proposes a multimodal ML framework fusing RGB, thermal,
and LiDAR data for pitting corrosion detection on offshore
infrastructure. The specific research objectives are:

\begin{enumerate}
    \item To design a hybrid feature-level fusion architecture
    integrating RGB, thermal infrared, and LiDAR modalities for
    pitting corrosion detection on offshore structural steel.
    \item To curate or synthesise a co-registered multimodal corrosion
    dataset representative of offshore environmental conditions,
    including surface fouling, variable illumination, and geometric
    irregularity.
    \item To evaluate the proposed framework against single-modality
    (RGB-only, IR-only, LiDAR-only) and dual-modality baselines
    using standardised pixel-level and detection-level metrics.
    \item To assess the robustness and generalisability of the
    framework under controlled simulations of realistic offshore
    degradation scenarios.
\end{enumerate}

The remainder of this paper is organised as follows:
Section~\ref{sec:related} reviews related work.
Section~\ref{sec:method} presents the proposed methodology.
Section~\ref{sec:outcomes} discusses expected outcomes.
Section~\ref{sec:timeline} outlines the research timeline.
Section~\ref{sec:conclusion} concludes the paper.

% ---------------------------------------------------------
\section{Related Work}
\label{sec:related}

\subsection{RGB-Based Corrosion Detection}

Visible-spectrum imaging is the most widely studied modality for
automated corrosion detection. Early work employed handcrafted
colour and texture features \cite{ref6}, whilst more recent studies
have leveraged deep CNNs, including YOLO-based detectors and
encoder-decoder segmentation networks, achieving F1-scores exceeding
0.90 on curated laboratory benchmarks \cite{ref7}. Transfer learning
from large-scale visual datasets has partially improved generalisation,
but RGB-based models remain vulnerable to illumination variation,
specular reflectance, and surface contamination—all of which are
endemic in offshore environments \cite{ref8}. Critically, RGB imaging
provides no information about subsurface material loss or surface
topography, limiting its diagnostic capability for pitting corrosion.

\subsection{Thermal Infrared-Based Detection}

Infrared thermography is an established NDT technique in which
subsurface defects—including corrosion-induced wall thinning—are
revealed as thermal anomalies under passive or active thermal
excitation \cite{ref9}. ML methods applied to thermal imagery,
including SVMs, CNNs, and U-Net segmentation architectures, have
demonstrated feasibility for corrosion characterisation \cite{ref10}.
Thermal imaging offers reduced sensitivity to ambient lighting
relative to RGB, enabling inspection in low-visibility or
night-time conditions. However, its spatial resolution is generally
inferior to optical cameras, and discriminating between
corrosion-induced and environmentally induced thermal gradients
remains an open problem in offshore deployments \cite{ref11}.

\subsection{LiDAR-Based Structural Inspection}

LiDAR sensors generate dense 3D point clouds with high geometric
fidelity, enabling pitting to be detected as localised surface
depressions \cite{ref12}. Deep learning approaches for point cloud
processing—including PointNet, PointNet++, and sparse 3D
CNNs—have been applied to structural inspection with promising
results \cite{ref13}. In marine and offshore contexts, LiDAR has
been used for hull inspection and pipeline monitoring, demonstrating
capability to resolve geometric features invisible to optical sensors
\cite{ref14}. Known challenges include sensitivity to water film,
maritime aerosol scatter, and specular reflectance from metallic
surfaces, all of which introduce noise in offshore point cloud
acquisitions.

\subsection{Multimodal Fusion Approaches and the Research Gap}

Multimodal fusion has consistently outperformed single-modality
baselines in domains including autonomous driving \cite{ref15},
aerial remote sensing \cite{ref16}, and medical image analysis
\cite{ref17}. Within structural inspection, limited studies have
fused RGB with thermal data for surface defect detection \cite{ref18}
or RGB with LiDAR for damage assessment \cite{ref19}. To the best
of the author's knowledge, no existing study has proposed a unified
framework fusing all three modalities—RGB, thermal infrared, and
LiDAR—specifically for pitting corrosion detection on offshore
infrastructure. This gap motivates the present research.

% ---------------------------------------------------------
\section{Proposed Methodology}
\label{sec:method}

\subsection{Research Design}

This study follows a quantitative experimental design. The primary
instrument is the proposed multimodal ML framework, evaluated through
controlled experiments on a purpose-assembled dataset. Ablation
studies isolating each modality's contribution, combined with
baseline comparisons, provide the empirical foundation for evaluating
the fusion approach.

\subsection{Data Acquisition and Preparation}

A multimodal dataset will be compiled from three complementary sources:

\begin{enumerate}
    \item \textbf{Public benchmarks:} Existing corrosion and surface
    defect datasets (e.g., CODEBRIM \cite{ref20}; NEU surface defect
    database \cite{ref21}) will be used for pre-training and transfer
    learning initialisation. Where paired thermal or 3D scan data are
    available, they will be incorporated.
    \item \textbf{Synthetic data generation:} Pitting textures, thermal
    anomaly maps, and point-cloud surface deformations will be
    synthesised using physically-based rendering pipelines and
    corrosion simulation models. Synthetic samples extend training
    coverage and allow controlled variation of pit depth, density,
    and environmental conditions.
    \item \textbf{Laboratory specimen imaging:} Corroded steel coupons
    representative of offshore structural grades will be imaged
    using an RGB camera, an uncooled microbolometer thermal camera,
    and a structured-light or time-of-flight LiDAR sensor under
    controlled conditions. Specimens will span pitted and non-pitted
    cases to support binary and severity-graded labelling.
\end{enumerate}

All acquisitions will be spatially co-registered to a common
reference frame via target-based multi-sensor calibration.
Ground-truth annotations will mark pit locations, pit severity
categories (shallow, moderate, severe), and corrosion-free
background regions, produced through a combination of automated
thresholding and manual expert labelling.

\subsection{Fusion Architecture}

The proposed framework adopts a \textbf{hybrid fusion strategy},
integrating feature-level cross-modal fusion with a late-stage
decision refinement mechanism, as illustrated in
Fig.~\ref{fig:framework}. Three parallel encoder branches process
the respective modalities:

\begin{itemize}
    \item \textbf{RGB branch:} A pre-trained CNN backbone
    (e.g., ResNet-50 or EfficientNet-B4) extracts hierarchical
    visual and texture features.
    \item \textbf{Thermal branch:} An analogous CNN backbone
    processes single-channel thermal images, extracting
    thermal gradient and anomaly maps.
    \item \textbf{LiDAR branch:} A PointNet++ or sparse 3D CNN
    encoder processes voxelised or range-projected LiDAR data,
    extracting surface curvature and depth deviation features.
\end{itemize}

Intermediate feature maps from the three branches are concatenated
at a fusion layer and passed through a shared decoder
(Feature Pyramid Network or U-Net-style architecture) for
pixel-level corrosion segmentation. A cross-modal attention
module will be investigated to learn adaptive, context-dependent
fusion weights, enabling the model to dynamically prioritise
modalities based on local signal quality.

\begin{figure}[htbp]
    \centering
    \fbox{%
        \parbox{0.92\linewidth}{%
            \centering\vspace{2.8cm}%
            \textit{[Framework Pipeline Diagram --- to be inserted]}%
            \vspace{2.8cm}%
        }%
    }
    \caption{Proposed trimodal fusion pipeline for pitting corrosion
    detection. Parallel CNN/PointNet++ encoders process RGB, thermal
    infrared, and LiDAR inputs. Features are fused via cross-modal
    attention and decoded to produce a pixel-level corrosion
    segmentation map with pit severity classification.}
    \label{fig:framework}
\end{figure}

\subsection{Evaluation Plan}

Model performance will be assessed using the following metrics:

\begin{itemize}
    \item \textit{Pixel-level:} Intersection over Union (IoU),
    Dice coefficient.
    \item \textit{Detection-level:} Precision, Recall, F1-score.
    \item \textit{Localisation:} Mean Average Precision (mAP) at
    IoU thresholds of 0.50 and 0.75.
\end{itemize}

Ablation experiments will systematically remove one or two modalities
to isolate each sensor's contribution to overall performance.
Robustness evaluation will introduce controlled degradations—simulated
sensor noise, illumination variation, partial surface occlusion—to
assess performance under realistic offshore conditions. Statistical
significance of performance differences between models will be
assessed using paired $t$-tests or Wilcoxon signed-rank tests
across cross-validation folds.

% ---------------------------------------------------------
\section{Expected Outcomes and Significance}
\label{sec:outcomes}

This research is expected to deliver the following contributions:

\begin{enumerate}
    \item \textbf{Novel trimodal fusion framework:} A first-of-its-kind
    ML pipeline unifying RGB, thermal, and LiDAR data for offshore
    pitting corrosion detection, with a modular architecture
    extensible to other structural inspection domains.
    \item \textbf{Benchmark experimental comparison:} Quantified
    performance comparisons between the proposed framework and
    single- and dual-modality baselines, providing the first
    empirical evidence for the value of trimodal fusion in this
    specific application.
    \item \textbf{Multimodal corrosion dataset:} A curated, co-registered
    collection of RGB, thermal, and LiDAR observations of corroded
    steel specimens, intended for public release to address current
    data scarcity in the field.
    \item \textbf{Practical deployment guidance:} Analysis of
    trade-offs between detection performance, sensor hardware cost,
    and computational overhead, informing feasibility assessments
    for real-world offshore inspection system design.
\end{enumerate}

Offshore structural inspection currently costs the industry several
billion US dollars annually, and undetected corrosion remains a
primary cause of structural incidents \cite{ref22}. An automated,
high-accuracy multimodal inspection framework could reduce reliance
on costly and hazardous manual inspection, support higher inspection
cadence, and provide earlier warning of critical corrosion
progression—directly advancing both structural safety and operational
efficiency in the offshore oil and gas sector.

% ---------------------------------------------------------
\section{Research Timeline}
\label{sec:timeline}

The proposed research spans approximately 24 months from
May~2026 to April~2028, organised into six sequential phases
as summarised in Table~\ref{tab:timeline}. Phases may overlap
where parallel tasks are feasible. Intermediate publication of
results is planned during Phase~5, targeting IEEE conferences or
journals in the areas of non-destructive testing, structural health
monitoring, or computer vision applied to engineering.

\begin{table}[htbp]
\caption{Research Timeline (May 2026 -- April 2028)}
\label{tab:timeline}
\centering
\renewcommand{\arraystretch}{1.25}
\begin{tabular}{p{0.52\linewidth} p{0.38\linewidth}}
\toprule
\textbf{Phase} & \textbf{Period} \\
\midrule
1. Literature Review \& Review Paper    & May -- Jul 2026   \\
2. Research Proposal                    & Jul -- Sep 2026   \\
3. Programming \& Tooling Setup         & Sep -- Nov 2026   \\
4. Dataset Acquisition \& Preprocessing & Nov 2026 -- Mar 2027 \\
5. ML Model Development \& Evaluation  & Mar -- Dec 2027   \\
6. Thesis Writing \& Dissemination     & Jan -- Apr 2028   \\
\bottomrule
\end{tabular}
\end{table}

% ---------------------------------------------------------
\section{Conclusion}
\label{sec:conclusion}

This paper has presented a research proposal for a multimodal ML
framework for pitting corrosion detection in offshore infrastructure,
integrating RGB, thermal infrared, and LiDAR sensor data through a
hybrid feature-level fusion architecture. The proposed work addresses
a documented gap in the literature—no prior study has unified all
three modalities for this application—and is expected to deliver a
novel, validated detection framework, a multimodal benchmark dataset,
and practical guidance for offshore structural health monitoring
system design. Successful completion of this research has the
potential to meaningfully advance inspection safety, reduce operational
costs, and improve the early detection of critical corrosion events
in the offshore oil and gas industry.

% ---------------------------------------------------------
% REFERENCES
% Replace each \bibitem with your actual source details.
% =========================================================
\begin{thebibliography}{99}

\bibitem{ref1}
[PLACEHOLDER] Author(s), ``Title,''
\textit{Journal/Conference}, vol.~X, no.~Y, pp.~ZZ--ZZ, Year.

\bibitem{ref2}
[PLACEHOLDER] Author(s), ``Title,''
\textit{Journal/Conference}, vol.~X, no.~Y, pp.~ZZ--ZZ, Year.

\bibitem{ref3}
[PLACEHOLDER] Author(s), ``Title,''
\textit{Journal/Conference}, vol.~X, no.~Y, pp.~ZZ--ZZ, Year.

\bibitem{ref4}
[PLACEHOLDER] Author(s), ``Title,''
\textit{Journal/Conference}, vol.~X, no.~Y, pp.~ZZ--ZZ, Year.

\bibitem{ref5}
[PLACEHOLDER] Author(s), ``Title,''
\textit{Journal/Conference}, vol.~X, no.~Y, pp.~ZZ--ZZ, Year.

\bibitem{ref6}
[PLACEHOLDER] Author(s), ``Title,''
\textit{Journal/Conference}, vol.~X, no.~Y, pp.~ZZ--ZZ, Year.

\bibitem{ref7}
[PLACEHOLDER] Author(s), ``Title,''
\textit{Journal/Conference}, vol.~X, no.~Y, pp.~ZZ--ZZ, Year.

\bibitem{ref8}
[PLACEHOLDER] Author(s), ``Title,''
\textit{Journal/Conference}, vol.~X, no.~Y, pp.~ZZ--ZZ, Year.

\bibitem{ref9}
[PLACEHOLDER] Author(s), ``Title,''
\textit{Journal/Conference}, vol.~X, no.~Y, pp.~ZZ--ZZ, Year.

\bibitem{ref10}
[PLACEHOLDER] Author(s), ``Title,''
\textit{Journal/Conference}, vol.~X, no.~Y, pp.~ZZ--ZZ, Year.

\bibitem{ref11}
[PLACEHOLDER] Author(s), ``Title,''
\textit{Journal/Conference}, vol.~X, no.~Y, pp.~ZZ--ZZ, Year.

\bibitem{ref12}
[PLACEHOLDER] Author(s), ``Title,''
\textit{Journal/Conference}, vol.~X, no.~Y, pp.~ZZ--ZZ, Year.

\bibitem{ref13}
[PLACEHOLDER] Author(s), ``Title,''
\textit{Journal/Conference}, vol.~X, no.~Y, pp.~ZZ--ZZ, Year.

\bibitem{ref14}
[PLACEHOLDER] Author(s), ``Title,''
\textit{Journal/Conference}, vol.~X, no.~Y, pp.~ZZ--ZZ, Year.

\bibitem{ref15}
[PLACEHOLDER] Author(s), ``Title,''
\textit{Journal/Conference}, vol.~X, no.~Y, pp.~ZZ--ZZ, Year.

\bibitem{ref16}
[PLACEHOLDER] Author(s), ``Title,''
\textit{Journal/Conference}, vol.~X, no.~Y, pp.~ZZ--ZZ, Year.

\bibitem{ref17}
[PLACEHOLDER] Author(s), ``Title,''
\textit{Journal/Conference}, vol.~X, no.~Y, pp.~ZZ--ZZ, Year.

\bibitem{ref18}
[PLACEHOLDER] Author(s), ``Title,''
\textit{Journal/Conference}, vol.~X, no.~Y, pp.~ZZ--ZZ, Year.

\bibitem{ref19}
[PLACEHOLDER] Author(s), ``Title,''
\textit{Journal/Conference}, vol.~X, no.~Y, pp.~ZZ--ZZ, Year.

\bibitem{ref20}
[PLACEHOLDER] Author(s), ``Title,''
\textit{Journal/Conference}, vol.~X, no.~Y, pp.~ZZ--ZZ, Year.

\bibitem{ref21}
[PLACEHOLDER] Author(s), ``Title,''
\textit{Journal/Conference}, vol.~X, no.~Y, pp.~ZZ--ZZ, Year.

\bibitem{ref22}
[PLACEHOLDER] Author(s), ``Title,''
\textit{Journal/Conference}, vol.~X, no.~Y, pp.~ZZ--ZZ, Year.

\end{thebibliography}

\end{document}
